\scanpage{183}

and which is trivial on $\widehat{\Pi}_{0,3}$, hence factors as
\[
(7)\qquad\widetilde N_{0,3}\longrightarrow\unsure{\mathfrak S_3|\mathcal G_3}.
\]
The kernels of (5), (7), are denoted $\widetilde M_{0,3}^{!}$, $\widetilde N_{0,3}^{!}$ respectively, so that we have a diagram of exact sequences

\unsure{\diagram{
\begin{tikzcd}[ampersand replacement=\&]
 \& \& \inserted{1}\arrow[d] \& \& \& \\
 \& \& \inserted{\Pi_{0,3}}\arrow[d] \& \& \& \\
(8) \& 1\arrow[r] \& \widetilde M_{0,3}^{!}\arrow[r]\arrow[d] \& \widetilde M_{0,3}\arrow[r]\arrow[d] \& \unsure{\mathfrak S_3|\mathcal G_3}\arrow[r]\arrow[d,equal] \& 1 \\
(9) \& 1\arrow[r] \& \widetilde N_{0,3}^{!}\arrow[r]\arrow[d] \& \widetilde N_{0,3}\arrow[r] \& \unsure{\mathfrak S_3|\mathcal G_3}\arrow[r] \& 1, \\
 \& \& \inserted{1} \& \& \&
\end{tikzcd}
}}

(9) being interpretable as the quotient of (8) by $\widehat{\Pi}_{0,3}$.

\inserted{\unsure{Let us note} also the \unsure{corresponding} extension structure
\[
(9\mathrm{bis})\qquad1\longrightarrow\widehat{\Pi}_{0,3}\longrightarrow\widetilde M_{0,3}^{!}\longrightarrow\widetilde N_{0,3}^{!}\longrightarrow1
\]}

We also have a canonical \unsure{homomorphism}
\[
(10)\qquad\unsure{\mathfrak S_3|\mathcal G_3}\simeq\widehat{\mathcal G}_{0,3}^{+}/\widehat{\Pi}_{0,3}\hookrightarrow\widetilde N_{0,3},
\]
whose composite with 7 is the identity, whence a canonical isomorphism
\[
(11)\qquad\widetilde N_{0,3}\simeq\unsure{\mathfrak S_3|\mathcal G_3}\times\widetilde N_{0,3}^{!}
\]
On the other hand, we also have an inclusion
\[
(12)\qquad
\underset{\substack{\unsure{\wr|\simeq}\\\unsure{\mathfrak S_3|\mathcal G_3}\times\mathbb Z/2\mathbb Z}}{\mathcal T_{0,3}}
\hookrightarrow\widetilde N_{0,3}
\]
which on the factor $\unsure{\mathfrak S_3|\mathcal G_3}\inserted{=\mathcal T_{0,3}^{+}}$ \unsure{coincides} with (10), and which \deleted{on} maps the nontrivial element $\tau$ of $\mathbb Z/2\mathbb Z\inserted{=\mathcal T_{0,3}^{!}}$ to the element, also denoted $\tau$, of the second factor $\widetilde N_{0,3}^{!}$ of (11).

\scanpage{184}
\mspage{391}

The composite homomorphism
\[
(13)\qquad \widetilde M_{0,3}\unsure{\hookrightarrow|\subset}
\unsure{\widehat{\widehat{\mathcal G}}_{0,3}}
\xrightarrow[\chi]{\text{\unsure{multiplier}}}
\widehat{\mathbb Z}^{*}
\]
\deleted{\unsure{gives|defines}} a homomorphism
\[
(14)\qquad \widetilde M_{0,3}\xrightarrow{\chi}\widehat{\mathbb Z}^{*}
\]
surjective, factoring through $\widetilde N_{0,3}$ and even \unsure{through} $\widetilde N_{0,3}/\unsure{\mathfrak S_3|\mathcal G_3}$

\diagram{\[
(15)\qquad
\begin{tikzcd}[ampersand replacement=\&]
\widetilde N_{0,3}\arrow[rr,"\chi"]\arrow[dr] \&\& \widehat{\mathbb Z}^{*} \\
\& \begin{gathered}\widetilde N_{0,3}/\unsure{\mathfrak S_3|\mathcal G_3}\\ \simeq\widetilde N_{0,3}^{!}\end{gathered}\arrow[ur] \&
\end{tikzcd}
,\]}

and giving rise to extension structures

\diagram{\[
(16)\qquad
\begin{tikzcd}[ampersand replacement=\&]
1\arrow[r] \& \widetilde M_{0,3}^{+}\arrow[r]\arrow[d] \& \widetilde M_{0,3}\arrow[r,"\chi"]\arrow[d] \& \widehat{\mathbb Z}^{*}\arrow[r]\arrow[d,equal] \& 1 \\
1\arrow[r] \& \widetilde N_{0,3}^{+}\arrow[r] \& \widetilde N_{0,3}\arrow[r,"\chi"] \& \widehat{\mathbb Z}^{*}\arrow[r] \& 1
\end{tikzcd}
,\]}

where we have set
\[
(17)\qquad
\left\{\begin{aligned}
\widetilde M_{0,3}^{+}&=\operatorname{Ker}(\widetilde M_{0,3}\to\widehat{\mathbb Z}^{*})\\
\widetilde N_{0,3}^{+}&=\operatorname{Ker}(\widetilde N_{0,3}\to\widehat{\mathbb Z}^{*})
\end{aligned}\right.
\]
-- the second row of (16) may be interpreted as the quotient of the first row by $\widehat\Pi_{0,3}$.

We have a canonical embedding
\[
(18)\qquad
\Gamma_{\mathbb Q}\overset{\mathrm{dfn}}{=}
\operatorname{Gal}\left(
\underbrace{\overline{\mathbb Q}_0}_{\inserted{\substack{\text{algebraic}\\\text{closure of }\mathbb Q\\\text{in }\mathbb C}}}
/\mathbb Q_0\right)
\unsure{\hookrightarrow|\subset}\widetilde N_{0,3}^{!}
\]
sending the \unsure{automorphism} of complex

\scanpage{185}

conjugation $\tau$ to the element already denoted $(\tau)$ (cf. (12)). The inverse image of \unsure{this|the} subgroup $\Gamma_{\mathbb Q}$ \deleted{in $\unsure{M}$}\inserted{of $\widetilde N_{0,3}^{!}$} under $\widetilde M_{0,3}\to\widetilde N_{0,3}^{!}\simeq\widetilde N_{0,3}/\unsure{\mathfrak S_3|\mathcal G_3}$, i.e. the inverse image of \deleted{$\unsure{\mathfrak S_3|\mathcal G_3}\times\Gamma_{\mathbb Q}\subset$}
\[
(19)\qquad
N_{0,3}\overset{\mathrm{dfn}}{=}
\unsure{\mathfrak S_3|\mathcal G_3}\times\Gamma_{\mathbb Q}
\subset\widetilde N_{0,3}
=\unsure{\mathfrak S_3|\mathcal G_3}\times\widetilde N_{0,3}^{!},
\]
is the extension \deleted{\unsure{of}} denoted $M_{0,3}$ elsewhere:

\margin{[We have a \deleted{\illegible{1}}\inserted{sub-}extension of \inserted{\unsure{the extension}} $\widetilde M_{0,3}$ of (2)]}

\diagram{\[
(20)\qquad
\begin{tikzcd}[ampersand replacement=\&]
1\arrow[r] \& \widehat\Pi_{0,3}\arrow[r]\arrow[d,"\unsure{\iota}"] \& M_{0,3}\arrow[r]\arrow[d,shift right=1.5pt]\arrow[d,no head,shift left=1.5pt,shorten <=7pt,shorten >=7pt] \& \overbrace{\unsure{\mathfrak S_3|\mathcal G_3}\times\Gamma_{\mathbb Q}}^{N_{0,3}}\arrow[r]\arrow[d,"{(\alpha,\beta)\to\beta}"] \& 1 \\
1\arrow[r] \& \widehat{\mathcal G}_{0,3}^{+}\arrow[r] \& M_{0,3}\arrow[r] \& \Gamma_{\mathbb Q}\arrow[r] \& 1
\end{tikzcd}
.\]}

Of course, it is in $M_{0,3}$ \inserted{and $N_{0,3}$} rather than in $\widetilde M_{0,3}$ \inserted{and $\widetilde N_{0,3}$} that we are interested -- hence in the remarkable subgroup $\Gamma_{\mathbb Q}$ of $\widetilde N_{0,3}^{!}$. We are going to ``determine'' by computation the group $\widetilde N_{0,3}^{!}$ -- and it will be a matter in the rest of the work of finding a maximum of necessary conditions on elements of $\widetilde N_{0,3}^{!}$ for them to come from $\Gamma_{\mathbb Q}$ -- \unsure{until} finding, if possible, a characterization (at least plausible, if not proved!) of the elements of $\Gamma_{\mathbb Q}$ -- i.e. of $\Gamma_{\mathbb Q}$ itself.

\scanpage{186}
\mspage{392}

II \uline{Computational determination of the elements of} $\widetilde N_{0,3}^{!}$.

1) \uline{Determination by} $h,g,f$.

\deleted{The \unsure{specification} of $u\in\unsure{\widetilde M_{0,3}^{!}|\widetilde N_{0,3}^{!}}$}

Recall that
\[
\begin{aligned}
(21)\qquad
&\unsure{\widetilde N_{0,3}^{!}|\widetilde M_{0,3}^{!}}
=\text{automorphisms of }\widehat{\mathcal G}_{0,3}^{+}
\simeq\unsure{\mathrm{Sl}(2,\mathbb Z)^{\wedge\prime}|\widehat{\mathrm{Sl}(2,\mathbb Z)^{\prime}}}
\text{ that stabilize}\\
&\widehat\Pi_{0,3},\quad\text{induce the identity on }
\unsure{\mathfrak S_3|\mathcal G_3}
=\widehat{\mathcal G}_{0,3}^{+}/\widehat\Pi_{0,3},\\
&\text{and preserve the loop structure (defined}\\
&\text{by the class of }\varepsilon_0\text{)}
\end{aligned}
\]

A $u\in\widetilde M_{0,3}^{!}$ is determined by a system
\[
(22)\qquad (h,g,f)\quad
\deleted{\unsure{\in(\widehat{\Pi}_{0,3})^{3}}\,\unsure{\times}\,\unsure{\widehat{\Pi}_{0,3}}}
\qquad(h,g,f\in\widehat\Pi_{0,3})
\]
\margin{$f$ determined modulo $f\mapsto\ell_0^{n}f$, $n\in\widehat{\mathbb Z}$}
\deleted{\unsure{satisfying}}\inserted{via} the formulas
\[
(23)\qquad
\left\{\begin{aligned}
u(\sigma_\infty)&=h\sigma_\infty\\
u(\rho)&=g\rho
\end{aligned}\right.
\]
\deleted{$\unsure{u(\varepsilon_0)}=\unsure{\operatorname{int}}\unsure{(}$}

the condition that these formulas define an automorphism (or at least an endomorphism) of $\widehat{\mathcal G}_{0,3}^{+}$ (which will necessarily induce the identity on $\unsure{\mathfrak S_3|\mathcal G_3}$) is expressed by the formulas
\[
(24)\qquad (h\sigma_\infty)^{2}=(g\rho)^{3}=1,
\]
which can also be written
\[
(25)\qquad
\boxed{\left\{\begin{aligned}
h\sigma_\infty(h)&=1\\
g\rho(g)\rho^{\unsure{2}}(g)&=1
\end{aligned}\right.},
\]
and the condition that $u$ preserves the loop structure is expressed by
\[
(26)\qquad u(\varepsilon_0)=\operatorname{int}(f^{-1})(\varepsilon_0^{\mu})
\qquad \mu\in\widehat{\mathbb Z}^{*}\text{ suitable},
\]
\margin{loop relation in $\widehat{\mathcal G}_{0,3}^{+}$ (raw form)}
\unsure{equivalent|equivalent} to the condition

\scanpage{187}

\margin{loop relation in $\widehat\Pi_{0,3}$ (raw form)}
\[
(27)\qquad
u(\ell_0)=\operatorname{int}(f^{-1})(\ell_0^{\mu}),
\]
it is easily made explicit in terms of $h,g,f$ (preferably in form (26)) taking into account that $\varepsilon_0=\sigma_\infty\rho$
\[
\underbrace{(h\sigma_\infty)(g\rho)}_{h\sigma_\infty(g)\underbrace{\sigma_\infty\rho}_{\varepsilon_0}}
=\underbrace{\operatorname{int}(f^{-1})\varepsilon_0^{\mu}}_{f^{-1}\varepsilon_0^{\mu}f},\quad\text{i.e}
\]
\[
f h\sigma_\infty(g)\underbrace{\varepsilon_0 f^{-1}}_{\unsure{\varepsilon_0(f^{-1})}\varepsilon_0=\sigma_\infty(\rho(f)^{-1})\varepsilon_0}
=\varepsilon_0^{\mu}\qquad\text{i.e.}
\]
\[
f h\sigma_\infty(g\rho(f)^{-1})=\varepsilon_0^{\mu-1}
\qquad\text{i.e., setting}
\]
\[
(28)\qquad \unsure{\nu}\overset{\mathrm{dfn}}{=}\frac{\mu-1}{2}\in\widehat{\mathbb Z}
\]
\margin{possible, since $\mu\in\widehat{\mathbb Z}^{*}$ hence $\mu\equiv1\ (2)$}
the relation (taking into account that $\varepsilon_0^{2}=\ell_0$)
\[
f h\sigma_\infty(g\rho(f)^{-1})=\ell_0^{\nu},
\]
or again, applying $\sigma_\infty$, since $\sigma_\infty(\ell_0)=\ell_1$
\[
\sigma_\infty(f)\sigma_\infty(h)g\rho(f)^{-1}=\ell_1^{\nu}
\]
that is, finally
\[
(29)\qquad
\boxed{\sigma_\infty(h)g=\sigma_\infty(f)^{-1}\ell_1^{\nu}\rho(f)}
\]
(\uline{loop relation in} $h,g,f$).

\margin{\uline{N.B.} $\nu$\inserted{$\in\widehat{\mathbb Z}$} is uniquely determined by $g,h$, and \unsure{we have} also $2\nu+1\in\widehat{\mathbb Z}^{*}$ $2\nu+1$ is the multiplier\ldots}

\uline{N.B.} One must in addition require, for $h,g,f$ to really define an \uline{automorphism}, that (24) be a \uline{presentation-generation} of $\widehat{\mathcal G}_{0,3}^{+}$ -- in what follows, we shall always tacitly take this ``preliminary condition'' as understood.